\section{The EpiSpace Data Engine}
\label{sec:engine}

\subsection{Acquisition contract and raw-channel isolation}

EpiSpace consumes existing OmniGibson acquisition bundles without mutating
them.  Each bundle contains a trajectory plan, ordered observations, camera
transforms, an entity inventory, RGB/depth/instance sensor arrays, and quality
reports.  The sweep plan is the authoritative job ledger.  The frozen policy
accepts only jobs marked \texttt{passed}; the adapter additionally requires all
core files, contiguous view indices, non-empty entities, integrity
\texttt{pass}, and a compatible trajectory class.  Review and failed jobs are
preserved in a reason-coded ledger.

Human previews concatenate RGB, depth, and instance-ID panels and are therefore
forbidden as model input.  The loader extracts the \texttt{uint8} RGB array from
each sensor archive into a release-owned PNG, records its dimensions and raw
pixel SHA-256, and stores instance pixel counts in a sidecar used only during
compilation.  A purity gate decodes every referenced image, verifies
1024$\times$1024 RGB content against the sidecar, and rejects paths under a
preview directory.  Target-view renders are checked against all SFT inputs by
both path and decoded RGB hash.

Across seven sweeps, \EpiPlannedJobs{} planned jobs produce
\EpiStrictBundles{} strict bundles.  Of these, \EpiSourceTrajectories{} are
write sources with \EpiSourceViews{} raw RGB views; 21 T10 bundles are verifier
sources only.  A fixed hash ranking assigns all descendants of each scene to
32/7/7 train/validation/test scenes.  No trajectory, sibling family, or target
render crosses that boundary.

\input{generated/dataset_table}

\subsection{Trajectory semantics and question legality}

Camera motion determines which supervision is identifiable.
\begin{itemize}
  \item \textbf{T1 walkthroughs} support overlap, revisit, separated-view
        registration, memory, metric, relation, and perspective reads.
  \item \textbf{T3 rotation stations} change yaw at a fixed position.  We
        supervise two-frame change detection only when exactly one eligible
        object enters view; metric tasks are rejected because rotation adds no
        translational baseline.
  \item \textbf{T4 adaptive arcs} keep a focus instance visible across a
        collision-free partial orbit.  The safe target is observed-instance
        identity, not canonical front or an unseen backside; these uniformly
        positive items are training-only.
  \item \textbf{T7 elevation ladders} change height and pitch at a station.
        Two sufficiently visible, unique objects are compared in an explicitly
        declared anchor-camera frame.
  \item \textbf{T8 reveal paths} provide a decisive category-presence view and
        matched evidence-deletion controls when such a family is valid.
  \item \textbf{T10 target views} never enter a write prefix.  Their RGB is an
        oracle-only observation used to grade visibility predicted from a
        same-scene write trajectory.
\end{itemize}
T2, T5, T6, and T9 remain expansion slots and are not counted as pilot assets.

Source questions must have an accepted/unknown status, a terminal source
certificate, and evidence views inside the actual model sequence.  Instance
questions are retained only when category-plus-view resolves one entity and at
least one declared anchor contributes 500 instance pixels.  T10 origin,
facing, and candidate entities require their own unique observed anchors.
Direction surfaces explicitly name the frame.  All 241 observed ontology
categories are deterministically localized into Chinese; an audit searches
model-visible text for residual raw labels.

\subsection{Geometry first, language second}

The compiler fixes an entity binding and structured target before realizing a
question.  Relations are computed in the declared camera frame; distances are
rounded only after metric evaluation; last-seen labels are recomputed over the
actual presented view list; and unknown means that the required category or
referent is absent from those inputs, not absent from simulator world truth.
Natural-language templates then verbalize the fixed target.  This ordering
prevents a data-generation paraphraser from defining label geometry or reading
the answer-side certificate; it does not prevent a trained student from making
geometric errors.  Future paraphrasers must consume only question-side facts
and remain answer-blind.

Each accepted item receives a globally unique fact ID, exact model/evidence
views, entity IDs in hidden metadata, a typed program, a Chinese surface
answer, and a replay certificate.  The independent replayer does not trust a
stored \texttt{passed} bit: it loads the source bundle and recomputes the target
for all 17 task/program types.  The frozen release replays
\EpiReplayedCertificates{}/\EpiReplayedCertificates{} targets successfully.

\subsection{Counterfactual family compiler}

The generic family compiler expands validated atomic facts rather than relying
only on rare hand-authored trajectories.
\begin{itemize}
  \item \textbf{Presence evidence (85 triples):} one background RGB is shared
        by all siblings; the second RGB is either a category-free distractor,
        a $\geq500$-pixel decisive view, or a different distractor.
  \item \textbf{Cross-view evidence (5 triples):} two RGBs, including the
        canonical anchor and one resolvable referent, are shared; only the
        third image reveals or withholds the second never-co-visible referent.
  \item \textbf{Trajectory reveal (1 triple):} a T8 decisive view is exchanged
        against a matched category-free view.
  \item \textbf{Frame equivariance (\EpiFrameFamilies{} pairs):} the same two
        images and objects are retained, but forward is anchored to image one
        versus image two; only pairs whose oracle relation changes are kept.
  \item \textbf{Claim verification (\EpiClaimFamilies{} pairs):} each certified
        cross-view relation yields balanced true and false claims with a
        structured corrected relation.
\end{itemize}
This produces \EpiConsistencyFamilies{} complete families.  Missing siblings,
unequal image counts, inconsistent family IDs, and cross-split membership are
release-blocking errors.

\subsection{Exports and fail-closed validation}

One intermediate episode IR is rendered into observe-first SFT, paired
isolated SFT, optional observable-state SFT, RLVR-ready prompts, and benchmark
rows.  The full diagnostic benchmark contains \EpiBenchmarkRecords{} records.
The recommended core is the union of all scene-disjoint family rows and all
semantic-signature holdouts: \EpiCoreBenchmarkRecords{} unique rows, with
separate \texttt{family} and \texttt{composition} views for disaggregated
reporting.

Eighteen release gates cover program typing; atom-seen/signature-unseen
composition; scene and family locks; complete sibling shapes; shared family
IDs; identical episode/isolated fact multisets; globally unique fact and record
IDs; target-view path and hash isolation; raw-RGB purity; functional
input--target dependence; terminal certificates; passing certificate checks;
and full independent replay.  All pass.  The manifest also retains
\EpiRejectedItems{} exclusions, dominated by ambiguous references, failed or
review acquisition jobs, trajectory-illegal tasks, small referents, and
verifier-only bundles.  Exclusions are part of the reported yield rather than
silently discarded data.
